% ==============================================================================
% ================================== EINLEITUNG ================================
% ==============================================================================

\chapter{Einleitung} 
\label{sec:einleitung}

\section{Hinführung und Relevanz des Themas}
\footnote{Aus Gründen der besseren Lesbarkeit wird in dieser Arbeit das generische Maskulinum verwendet. Weibliche und anderweitige Geschlechtsidentitäten sind ausdrücklich mitgemeint.}
Kiefer betont, dass die zwischenmenschliche Kommunikation in der Führungspraxis und in dienstleistungsorientierten Arbeitsbeziehungen das zentrale Werkzeug für die Interaktion sowie für die Ausgestaltung sozialer Beziehungen darstellt \cite[S.~1]{3}. Kommunikationsprozesse verlaufen jedoch selten vollkommen störungsfrei. Insbesondere in fehlerbehafteten Leistungssituationen entstehen Spannungen, die ohne ein adäquates Konfliktmanagement rasch eskalieren können. Ein professioneller Umgang mit kommunikativen Störungen erfordert von Dienstleistern sowohl theoretisches Orientierungswissen als auch Handlungskompetenz zur Deeskalation und kundenorientierten Konfliktlösung \cite[S.~1--4]{5}.

\section{Problemstellung und Zielsetzung}
Gegenstand der vorliegenden Arbeit ist die wissenschaftliche Analyse einer eskalierten Gesprächssituation zwischen einem Dienstleister im Bereich der Prozessautomatisierung (\ac{RPA}/\ac{KI}) und einer Kundin aus der Finanzabteilung. Ausgangspunkt ist ein operativer Fehler der Dienstleistungsseite, der im anschließenden persönlichen Klärungsgespräch eine Eskalationsdynamik auslöst: Emotionale Reaktionen, gegenseitige Vorwürfe und schließlich der Abbruch der Kommunikation durch die Kundin sind die Folge.

Das Ziel der Arbeit besteht darin, den Verlauf dieser Interaktion auf Basis etablierter kommunikationstheoretischer Modelle zu untersuchen, die Ursachen der Eskalation zu identifizieren und praxistaugliche Handlungsoptionen zur Wiederherstellung einer konstruktiven Kundenbeziehung zu entwickeln.

\section{Aufbau der Arbeit}
Die Ausarbeitung gliedert sich im Hauptteil entlang der spezifischen Aufgabenstellungen in folgende Abschnitte:

In Abschnitt 2 wird das persönliche Treffen in Form eines vollständigen Dialogskripts konstruiert. Dieses dient als Analysebasis für die nachfolgenden Untersuchungsschritte.

In Abschnitt 3 erfolgt die Erläuterung und Anwendung der fünf pragmatischen Axiome nach Watzlawick et al. auf konkrete Passagen des erstellten Dialogs, um die Wirkungsmechanismen der gestörten Kommunikation nachzuweisen.

In Abschnitt 4 werden operative und kommunikative Folgemaßnahmen erarbeitet, die der Dienstleister ergreifen muss, um die Kundin trotz des gescheiterten Erstgesprächs zu halten und ein zweites Gespräch vorzubereiten.

Abschnitt 5 beleuchtet weiterführende Handlungsalternativen für den Fall, dass auch das zweite Klärungsgespräch nicht zu einem Konsens führt.

In Abschnitt 6 wird der Gesprächsverlauf retrospektiv analysiert, um die kritischen Wendepunkte sowie die zugrundeliegenden Ursachen der Eskalation systematisch offenzulegen.

Abschnitt 7 fasst die wesentlichen Erkenntnisse im Rahmen einer abschließenden Reflexion zusammen und bewertet die kritischen Erfolgsfaktoren kundenorientierter Konfliktbewältigung.

% ==============================================================================
% ================================== EINLEITUNG ================================
% ==============================================================================